\documentclass[10pt]{article}

\pagestyle{plain}

\title{{\small 1999 IMACS Conference on Applications of Computer Algebra (IMACS-
ACA'99)}\\
Interaction on\\
Physique, Number Theory \& Computer Science}
\author{Session organized by\\
{\bf Hoang Ngoc Minh}\thanks{e-mail:hoang@lifl.fr\hfill http://www.lifl.fr/\~{}h
oang},\\
University de Lille II, 59624, Lille C\'edex,\\
\&\\
LIFL - URA 369 CNRS, 59655 Villeneuve d'Ascq, France.}
\date{}

\begin{document}

\maketitle

\section{Session Description}

The interaction on physics, number theory and computer science is the old one.
The recent progress of symbolic computation on Feymann diagrams,
knot invaviants and chord diagrams lead to calculation of transcendental numbers
(multiple zeta values, harmonic series, ...)
and need new algorithms on algebraic methods.
The software developments and their intensive experimentations lead to very
deeper conjectures on number theory.

The aim of this session is to present several examples of how computers can be
very useful in obtaining new results, new conjectures and ... new mathematics,
via the following topics

\section{Session Topics}
\begin{itemize}
\item ACA in High Energy Physics, Quantum Field Theory, Statistical Physics,
\item ACA in G-Functions, Hypergeometric Functions, Polylogarithms,
\item ACA in Harmonic Series, Multiple Zeta Values,
\item ACA in Hecke Algebra, Quasi-Hopf Algebra, Quasi-Symetric Functions,
\item $\ldots$
\end{itemize}

\end{document}
